\section{Scope and non-claims}\label{sec:scope}
% statements-of-record rows resolved here: none

The closure developed in this paper is deliberately conditional.  The
point of the construction is to isolate exactly which recognition
sources and imported theorem stacks would close the Strong-BSD scalar
identity inside the Six Birds closure discipline, and to record what
is not being claimed.  The following seven non-claims are part of the
statement of the paper.  They fix the intended reading of the theorem,
the role of the recognition sources, the status of the imports, and
the meaning of the final AOR membership assertion.

\begin{itemize}[leftmargin=*]
\item[\textbf{NC-1.}]
\textbf{Not unconditional BSD.}  This is not an unconditional
classical proof of BSD, scalar or strong.  The result is theorem-grade
within the Six Birds closure discipline, conditional on the standing
shell-closure record, the named recognition sources, and the headline
imported theorem stacks.  Outside the framework it reads as the
conditional theorem
\[
  \begin{aligned}
  &\SelBSD\text{-closure}
    \wedge\ \GammaPD \wedge \GammaSP \wedge \GammaGZ\\
  &\quad{}\wedge\ \chiCTp\text{-imports}
    \wedge\ \AofE\text{-imports}
    \wedge\ \text{Beilinson-imports}\\
  &\Longrightarrow \mathrm{Strong\ BSD}.
  \end{aligned}
\]
This is the headline reading of the paper.  The closure is a precise
conditional implication, not a replacement for the missing
unconditional arithmetic theorems.

\item[\textbf{NC-2.}]
\textbf{Recognition sources not derived.}  No recognition source is
derived from framework primitives.  Each \(\Gamma_{\mathrm{BSD}}^{*}\)
is closure content of the formed BSD layer: a typed carrier whose
mathematical payload is supplied as a hypothesis and then consumed by
the shell.  The formalization records these carriers as explicit
structure data, not as axioms and not as hidden constants.  Thus the
three sources are neither proved by the framework nor smuggled into
the conclusion; they are named, typed, and kept visible throughout the
argument.

\item[\textbf{NC-3.}]
\textbf{Rank-\(\geq 2\) Gross--Zagier remains open.}  We do not prove
the rank-\(\geq 2\) generalized Gross--Zagier identity.  The relevant
missing identity is the higher-rank comparison between the leading
\(L\)-value and a Neron--Tate regulator determinant of algebraic
cycles, with the necessary normalizing factors made explicit.  The
no-go result used in this paper is an audit result: in the audited
theorem-grade literature, the needed number-field identity is not
available in the required form.  This is not a non-existence claim.
It says only that the present literature does not supply the theorem
needed to remove the corresponding recognition source.

\item[\textbf{NC-4.}]
\textbf{Additive \(p\leq 3\) Iwasawa main conjecture remains open.}
We do not prove the additive \(p\leq 3\) Iwasawa main conjecture with
the explicit local correction factors and characteristic-ideal
identity required by the closure.  This is again an audit-result
no-go, not a claim that such a theorem cannot exist.  The paper treats
the missing small-additive-prime input as open arithmetic content and
does not derive it from the closure framework.  In particular, the
odd-additive recognition carriers do not by themselves discharge the
\(p=2\) part of this obstruction; any such discharge would have to be a
separate supplied record or a future imported theorem.

\item[\textbf{NC-5.}]
\textbf{Imports not eliminated.}  The \(\chiCTp\) comparison remains
conditional on its imported Cassels--Tate stack.  Likewise, the
\(\AofE\) Selmer-complex pairing substrate and the Beilinson
rank-\(\geq 2\) comparison enter as imported theorem stacks.  These
structures are sourced assumptions for this paper, not results proved
here.  The closure derives consequences from them; it does not
eliminate them, reconstruct them, or claim independent proofs of the
classical arithmetic geometry they package.

\item[\textbf{NC-6.}]
\textbf{Anchor-level evidence is not family-level.}  The
overconvergent \(\eta\)-formula sub-residual is verified at the
anchors \(49\mathrm{a}1@7\), \(36\mathrm{a}1@3\), and
\(121\mathrm{b}1@11\).  Those checks are anchor-level numerical
evidence.  They are not a uniform family-level statement over all
elliptic curves or all additive odd primes in the relevant regime.
The conditional formula in \Cref{sec:eta-formula} remains conditional
on its published-import stack and its named nonvanishing/primitivity
input.

\item[\textbf{NC-7.}]
\textbf{AOR is non-eliminative.}  The AOR-instance reading does not
close the open obligations.  It reclassifies the closure records: the
recognition discharges are bridged records, meaning the corresponding
open arithmetic content is carried as supplied recognition content
rather than eliminated.  The rank-\(\geq 2\) generalized
Gross--Zagier identity and the additive \(p\leq 3\) Iwasawa main
conjecture therefore remain open after the AOR reading.  Moreover,
only the local syntactic refinement-stable predicate is mechanized in
this paper, not the full Audited Operational Realisability
meta-theory.
\end{itemize}

These non-claims also govern the wording of the paper.  We do not say
that the paper proves BSD, proves the additive-prime Iwasawa main
conjecture, or proves the higher-rank Gross--Zagier identity.  The
paper lands Strong BSD as a conditional closure.  The standing shell-closure
record is assumed; recognition sources are supplied and carried as
hypotheses; headline imported theorem stacks are imported and cited; the
AOR reading records a non-eliminative
audit-closed membership statement under those declared records.
\begin{table}[tbp]
\centering
\footnotesize
\caption{Summary of the non-claims of \Cref{sec:scope}.}
\label{tab:nonclaims}
\begin{tabularx}{\textwidth}{@{}>{\raggedright\arraybackslash}p{0.05\textwidth}
>{\raggedright\arraybackslash}p{0.27\textwidth}
>{\raggedright\arraybackslash}X@{}}
\toprule
 & Boundary & Reading \\
\midrule
1 & Not a proof of BSD & Conditional translation; the scalar hypothesis has
the strength of the conclusion. \\
2 & Hypotheses not derived & Recognition hypotheses and comparison inputs
are assumed. \\
3 & No necessity & Normalized \(\GammaGZ\) alone gives the scalar identity
in rank \(\ge2\). \\
4 & No arithmetic formed layer & Intertwining readout and original
fixedness are assumed. \\
5 & No computability & Named computability fields are supplied, not
proved. \\
6 & Open inputs & Rank-\(\ge2\) leading terms; additive \(p\in\{2,3\}\)
main conjectures. \\
7 & \(\eta\)-formula & Conditional certificate only; no evidence
claimed. \\
8 & Register & Statuses under a declared refinement; no arithmetic
proved. \\
\bottomrule
\end{tabularx}
\end{table}

